\chapter{Deployment scripts and tuning} \label{cap:deployment}

The deployment process was automated around two complementary concerns: creating the cloud infrastructure in a repeatable way and updating the Kubernetes workloads with minimal manual intervention. The repository therefore separates infrastructure provisioning, image publication, application deployment, monitoring, and backup operations into scripts under the \texttt{scripts/cloud} directory and declarative manifests under \texttt{deploy/k8s}.

\section{Infrastructure bootstrap}

The cloud layer is managed with Terraform. Before it is executed, the Terraform initialization script creates, when necessary, a Google Cloud Storage bucket for the remote Terraform state and enables object versioning on it. The bucket name is derived from the active Google Cloud project, which keeps the Terraform backend consistent between machines.

The Terraform configuration provisions the main managed services used by the system: a PostgreSQL Cloud SQL instance, the individual microservice databases, an Artifact Registry Docker repository, Pub/Sub topics, a backup bucket, a VPC and subnet, and the GKE cluster. The cluster uses a custom node pool with \texttt{e2-standard-2} nodes and Workload Identity enabled, so microservice pods can access Google Cloud resources through a Kubernetes service account. This is also used by the backup CronJob and the Cloud SQL proxy sidecars.

\section{Image build and publication}

Container images are built and pushed by the \texttt{build\_and\_push.sh} script. The script iterates over all runtime services, including the gateway and other services like Redis. Each service image is built with Docker for the target platform (which defaults to \texttt{linux/amd64} that is the system Cloud Run containers must use), tagged with the Artifact Registry path, and pushed to the \texttt{vehicles} repository.

\section{Kubernetes deployment}

The workloads and other configurations are set as Kubernetes manifests and composed with Kustomize. The root \texttt{deploy/k8s/kustomization.yaml} applies the shared namespace, service account, Cloud SQL configuration, all microservice deployments and services, Redis, the backup CronJob and the ingress rule. To aid the Kubernetes deployment, a helper script was developed (\texttt{scripts/cloud/k8s.sh}) that wraps common operational actions:

\begin{itemize}
    \item \texttt{up}: Applies the Kubernetes manifests;
    \item \texttt{down}: Deletes the Kubernetes manifests;
    \item \texttt{status}: Shows the current pod and service status;
    \item \texttt{restart}: Restart all deployments (rollout restart).
\end{itemize}

The script detects the namespace from the manifest, selects the kubeconfig from CI or from the local project file, checks cluster connectivity, applies the manifests and optionally waits for rollouts to complete.

After applying manifests, the script restarts deployments so that pods pull the latest image tag from the Artifact Registry. Ingress can be managed through the same deployment flow. The ingress routes all \texttt{/api/} traffic to the gateway service, with explicit WebSocket paths for chat and auctions and longer proxy read/send timeouts to support persistent connections.

\section{Operational tuning}

Much of the performance tuning was done at the Kubernetes workload level. Each service defines CPU and memory requests and limits so the scheduler can place pods reliably on the small GKE node pool and so individual services cannot consume all cluster resources during stress tests. Lighter services like \texttt{user} use smaller requests, while listing or search reserve more CPU and memory because they perform heavier database queries or aggregation work. Services that connect to Cloud SQL include a Cloud SQL proxy sidecar with its own resource budget, preventing the proxy from competing invisibly with the application container. These values were tweaked after CloudyDay to better strengthen the infrastructure.

Readiness and liveness probes were added to the HTTP and gRPC-facing deployments. Readiness probes prevent Kubernetes services from sending traffic to pods before they have completed startup and dependency initialization. Liveness probes allow Kubernetes to restart pods that become unhealthy.

Horizontal Pod Autoscalers are configured for the stateless application services. Most services scale from one replica up to eight replicas using CPU and memory utilization targets of roughly 70\% and 80\%, respectively. The chat service uses slightly lower targets and a maximum of ten replicas because WebSocket traffic can keep pods busy even when request counts are not high. For services such as listing, search, seller, user, etc., the HPA behavior allows fast scale-up while applying a 180 second stabilization window when scaling-down (to avoid oscillation during short traffic spikes while still allowing the cluster to reduce cost after load decreases).

\section{Backup and monitoring support}

Database and dataset backup operations are also scripted. The manual \texttt{backup\_data.sh} script exports the Cloud SQL databases to Google Cloud Storage and uploads prepared dataset artifacts. There is also, in the cluster, a Kubernetes CronJob that runs daily at 02:00 UTC and exports the configured Cloud SQL databases to the backup bucket. The backup bucket has object versioning enabled and a lifecycle rule that moves older objects to Coldline storage after 30 days, reducing storage cost while keeping recovery points available.

Monitoring is deployed separately with the monitoring Kubernetes helper and use the manifests under \texttt{deploy/k8s/monitoring}. This includes Prometheus, Grafana, kube-state-metrics, services, RBAC rules, etc. Keeping monitoring in its own namespace and script allows restart, remove or inspect observability components without touching the application workloads. During tuning, the Prometheus and Grafana views were used to compare pod resource usage with the configured requests, limits and HPA thresholds after smoke and stress tests.